\documentclass[10pt,a4paper]{amsart}
\usepackage[margin=19mm]{geometry}
\usepackage{amsmath,amssymb,mathtools}
\usepackage[scr=rsfs,cal=euler]{mathalfa}
\usepackage{xfrac}
\usepackage[unicode,hidelinks,bookmarks=false]{hyperref}
\newcommand{\ie}{i.e.\ }
\newcommand{\cf}{cf.\ }
\newcommand{\resp}{respectively\ }
\newcommand{\Subsection}{\subsection}
\providecommand{\nobreakdash}{\nobreak\hbox{-}}
\setlength{\emergencystretch}{2em}
\multlinegap0pt
\usepackage[shortlabels]{enumitem}
\usepackage{mathtools}
\usepackage{amssymb}
\usepackage{xcolor}
\usepackage{tikz}
\newtheorem{assumption}{Assumption}
\newtheorem{lemma}{Lemma}
\usepackage{cleveref}
\crefname{assumption}{Assumption}{Assumptions}
\crefname{lemma}{Lemma}{Lemmas}
\newcommand{\I}{\mathrm{I}}
\newcommand{\pw}{\mathrm{pw}}
\newcommand{\s}{\mathrm{s}}
\title{Lower eigenvalue bounds with hybrid high-order methods}
\author{Ngoc Tien Tran}
\date{Source: arXiv:2406.06244v3 (CC BY 4.0)}
\begin{document}
\maketitle

\noindent\emph{Source and licence.} Lower eigenvalue bounds with hybrid high-order methods. Version arXiv:2406.06244v3 (CC BY 4.0), distributed by arXiv under the Creative Commons Attribution 4.0 International licence, http://creativecommons.org/licenses/by/4.0/, as recorded in the arXiv licence metadata for this exact version and shown on the abstract page https://arxiv.org/abs/2406.06244v3. Full licence notice: \texttt{licenses/arxiv-2406.06244-CC-BY-4.0.txt}.\par\smallskip

\noindent\textit{Editorial scope.} Counted unit: the article's lemma lem:linear-independency (source0.tex lines 360--369). The article's complete original proof of it is delivered with it (source0.tex lines 370--401, one \texttt{proof} environment of 32 lines). No mathematics is written, repaired or re-derived: the statement and the proof are the article's own lines, in the article's own notation, and no retained line is substituted or re-flowed.  Retained uncounted as the source-local context the proof needs: lines 179--181; lines 259--261; lines 299--301; lines 304--306; lines 308--310; lines 349--351.  Nothing was omitted from any retained span, so the package carries no omission record.  No retained line contains a citation, so the wrapper carries no bibliography and no bibliography entry.  No retained line contains a figure environment, an image-inclusion command or drawing code, so no image is delivered and the package declares no asset dependency.  Editorial formatting only, in addition: an AMS \texttt{amsart} preamble that restates the article's own declarations and macros used by the retained lines, namely the environments \texttt{assumption}, \texttt{lemma} on separate counters, together with the standard packages those lines need.  The retained environments print in this document as Assumption 1, Lemma 1 in order of appearance, whereas the article prints the same items with the numbering of its own full document.  The complete native source of the article as distributed in the arXiv e-print for this exact version is bundled unchanged as \texttt{sources/160547/source0.tex}.
\begin{align}\label{def:a-orthogonality}
	a(u(j), u(k)) = b(u(j), u(k)) = 0 \quad\text{for any } j \neq k.
\end{align}

\begin{align}\label{eq:min-max-principle-discrete}
	\lambda_h(j) = \min_{W_h \subset V_h, \mathrm{dim}(W_h) = j} \max_{v_h \in W_h \setminus\{0\}} (\|v_h\|_{a_h}^2 + \|v_h\|_{\s_h}^2)/\|v_h\|_{b_h}^2.
\end{align}

\begin{assumption}[$a$-orthogonality]\label{A}
	$\|\I_h v\|_{a_h}^2 \leq \|v\|^2_a - \|v - \mathrm{G}_h v\|_{a_\pw}^2$
\end{assumption}

\begin{align}\label{ineq:v-Gv<=v}
	\|v - \mathrm{G}_h v\|_{a_\pw} \leq \|v\|_a.
\end{align}

\begin{assumption}[scaling of stabilization]\label{B}
	$\|\I_h v\|_{\s_h}^2 \leq \alpha \|v - \mathrm{G}_h v\|_{a_\pw}^2$
\end{assumption}

\begin{assumption}[$b$-orthogonality]\label{C}
	$\|\I_h v\|_{b_h}^2 \geq \|v\|_b^2 - \beta \|v - \mathrm{G}_h v\|_{a_\pw}^2$.
\end{assumption}

\begin{lemma}[linear independency]\label{lem:linear-independency}
	If $\beta \lambda(j) < 1$, then (a)--(b) hold.
	\begin{enumerate}[wide]
		\item[(a)] The interpolations $\I_h u(1)$, $\dots$, $\I_h u(j)$ of the exact eigenvectors $u(1), \dots, u(j)$ are linearly independent in $V_h$.
		\item[(b)] Suppose additionally $\lambda_h(j) < \infty$, then there exists $v \in W$ with $\|v\|_b = 1$ and
		\begin{align*}
			(1 - \alpha - \beta \lambda_h(j))\|v - \mathrm{G}_h v\|_{a_\pw}^2 \leq \lambda(j) - \lambda_h(j).
		\end{align*}
	\end{enumerate}
\end{lemma}

\begin{proof}[Proof of (a)]
	Suppose otherwise, then there are real numbers $t_1, \dots, t_j$ such that $v \coloneqq t_1 u(1) + \dots + t_ju(j)$ satisfies $\|v\|_b = 1$ but $\I_h v = 0$.
	Elementary algebra with the orthogonality in \eqref{def:a-orthogonality} proves
	$1 = \|v\|_b^2 = t_1^2 \|u(1)\|_b^2 + \dots t_j^2 \|u(j)\|_b = t_1^2 + \dots + t_j^2$ and therefore,
	\begin{align}\label{ineq:bound-a-v-v}
		\|v\|_a^2 = t_1^2 \|u(1)\|^2_a + \dots + t_j^2 \|u(j)\|_a^2 &= t_1^2 \lambda(1) + \dots + t_j^2 \lambda(j)\nonumber\\
		&\leq (t_1^2 + \dots + t_j^2)\lambda(j) = \lambda(j).
	\end{align}
	The combination of this with \Cref{C} and \eqref{ineq:v-Gv<=v} results in
	\begin{align}\label{ineq:proof-dLEB-claim-1}
		1 = \|v\|_b^2 \leq \beta \|v - \mathrm{G}_h v\|_{a_\pw}^2 \leq \beta \|v\|^2_a \leq \beta \lambda(j);
	\end{align}
	a contradiction to $\beta\lambda(j) < 1$.\\[-0.5em]
	
	\paragraph{\emph{Proof of (b)}}
	Since the interpolations $\I_h u(1), \dots, \I_h u(j)$ are linearly independent from (a), $\I_h W \subset V_h$ is a $j$-dimensional subspace of $V_h$.
	The Rayleigh-Ritz principle \eqref{eq:min-max-principle-discrete} on the discrete level shows
	\begin{align*}
		\lambda_h(j) \leq \max_{v_h \in \I_h W\setminus\{0\}} (\|v_h\|_{a_h}^2 + \|v_h\|_{\s_h}^2)/\|v_h\|_{b_h}^2.
	\end{align*}
	Let $v \in W \setminus \{0\}$ with $\|v\|_b = 1$ be given such that $\I_h v$ is a maximizer of this Rayleigh quotient (which can attain the value $+\infty$).
	Then
	\begin{align}\label{ineq:proof-dLEB-min-max}
		\lambda_h(j) \|\I_h v\|_{b_h}^2 \leq \|\I_h v\|_{a_h}^2 + \|\I_h v\|_{\s_h}^2.
	\end{align}
	From \Cref{A} and $\|v\|_a^2 \leq \lambda(j)$ in \eqref{ineq:bound-a-v-v}, we deduce that
	$\|\I_h v\|_{a_h}^2 \leq \lambda(j) - \|v - \mathrm{G}_h v\|_{a_\pw}^2$. This and \Cref{B} imply that the right-hand side of \eqref{ineq:proof-dLEB-min-max} is bounded by $\lambda(j) - (1 - \alpha)\|v - \mathrm{G}_h v\|_{a_\pw}^2$, while $\lambda_h(j)(1 - \beta \|v - \mathrm{G}_h v\|_{a_\pw}^2)$ is a lower bound for the left-hand side of \eqref{ineq:proof-dLEB-min-max} thanks to \Cref{C} and $\|v\|_b = 1$. Hence,
	\begin{align*}
		\lambda_h(j)(1 - \beta \|v - \mathrm{G}_h v\|_{a_\pw}^2) \leq \lambda(j) - (1 - \alpha)\|v - \mathrm{G}_h v\|_{a_\pw}^2.
	\end{align*}
	Reorganizing these terms concludes the proof of (b).
\end{proof}

\end{document}
